Before reporting reduced-form leakage tests, we describe the magnitude of the EU ETS cost shock at two levels: from the perspective of the regulated firm itself (how large a fraction of its operating cost is attributable to allowance shortfalls priced at the prevailing EUA?), and from the perspective of the downstream buyer (how large a fraction of the buyer's input bill is supplied by ETS-regulated firms, and what cost-pass-through implies for the buyer's relative-price signal?). Throughout this section, all numbers are computed on the full Belgian panel.

Figure~\ref{fig:eua_price_timeseries} sets the aggregate context: the daily EUA spot price across the available ICAP window (2010--2022). Phase~II prices were depressed near \EUR{5}--\EUR{15}; the Phase~II/III transition (Jan 2013) coincided with a further price decline; the October 2015 MSR decision committed the EU to a rules-based mechanism for absorbing the allowance surplus, with realised prices rising from \EUR{6}--\EUR{8} in 2015--16 to above \EUR{25} in 2018--19 and above \EUR{80} in 2022 after the Phase IV cap entered force.

\begin{figure}[H]
    \centering
    \includegraphics[width=\textwidth]{phase3_eua_price_timeseries.pdf}
    \caption{Daily EUA spot price (ICAP primary market, EUR/tCO\textsubscript{2}), 2010--2022. Red solid lines mark the Phase~II/III (Jan 2013) and Phase~III/IV (Jan 2021) boundaries. Green dashed line marks the October 2015 MSR decision (Decision (EU) 2015/1814). ICAP daily series begins in 2010, mid-Phase~II; Phase~I (2005--07) and most of Phase~II (2008--09) are not covered.}
    \label{fig:eua_price_timeseries}
\end{figure}

\subsection{The shock to treated firms}\label{sec:shock_firms}

For each ETS-regulated firm-year we compute the same-year carbon-cost-share intensity as
\[
    \text{firm cost share}_{j,t} \;=\; \frac{\text{shortage}_{j,t} \cdot \text{EUA}_t}{\text{total cost}_{j,t}}
\]
where $\text{shortage}_{j,t} = \max(0, \text{emissions}_{j,t} - \text{free allowances}_{j,t})$, EUA$_t$ is the average annual allowance price, and total cost is from Annual Accounts. This is the descriptive analog of the time-invariant treatment intensity in equation~\eqref{eq:fcs}: it answers ``what fraction of this firm's operating cost would the firm have to pay for unallocated emission allowances at this year's EUA price?'' Distributions by phase are reported in Table~\ref{tab:firm_cost_share_by_phase}.

\begin{table}[H]
    \centering
    \caption{Distribution of carbon cost share across EU ETS Phases.}
    \label{tab:firm_cost_share_by_phase}
    \footnotesize
    \begin{tabular}{lccccccc}
        \toprule
        & & & & & & \multicolumn{2}{c}{Mean} \\
        \cmidrule(lr){7-8}
        Phase & N & Share shortage $> 0$ & p50 & p90 & p99 & Unweighted & Sales-weighted \\
        \midrule
        I (2005--07)            & 505 & 19.80\% & 0.00\% & 0.02\% & 0.46\% & 0.02\% & 0.08\% \\
        II (2008--12)           & 849 & 12.50\% & 0.00\% & 0.01\% & 0.46\% & 0.18\% & 0.08\% \\
        III pre-MSR (2013--17)  & 917 & 62.20\% & 0.01\% & 0.30\% & 8.67\% & 0.56\% & 0.09\% \\
        III post-MSR (2018--20) & 511 & 72.00\% & 0.09\% & 1.40\% & 24.00\% & 1.41\% & 0.31\% \\
        \textbf{IV (2021--22)}  & \textbf{319} & \textbf{83.40\%} & \textbf{0.34\%} & \textbf{2.60\%} & \textbf{59.90\%} & \textbf{6.85\%} & \textbf{0.71\%} \\
        \bottomrule
    \end{tabular}
\end{table}

The Phase-by-Phase pattern is sharp. Phase~I and Phase~II shocks are essentially zero everywhere because free allocation generously covered ETS firms' emissions; the median firm has zero shortage. Phase~III pre-MSR begins to bite in the upper tail (p99 = 8.7\%) but the median firm still has near-zero shortage. The MSR-driven price rise pushes Phase~III post-MSR p90 to $1.40\%$ and Phase~IV p90 to $2.60\%$, with the right tail reaching $\sim$60\% of total cost at the most-exposed installations.

A complementary metric --- the effective carbon price per tonne emitted --- shows the same pattern at the aggregate level. The effective price, defined as the ratio of total carbon cost paid (shortage $\times$ EUA, summed across firms) to total emissions, rises from $\$1.80$/tCO$_2$ in Phase~I to $\$32.82$/tCO$_2$ in Phase~IV, an 18-fold increase. Phase~IV is the only phase in which the carbon-pricing shock is unambiguously ``real'' at the population level.

\subsection{The shock to treated buyers}\label{sec:shock_buyers}

The magnitude of the shock at the buyer's perspective is what determines whether substitution incentives are operative. We compute three measures, each at a different denominator scale, on the panel of buyer $\times$ ETS-supplier $\times$ year cells.

\paragraph{Pair-shock at the NACE-4d denominator.} Letting $j$ denote an ETS supplier and $b$ a buyer, define
\[
    \text{pair-shock}_{j,b,t} \;=\; \text{firm cost share}_{j,t} \;\times\; \frac{\text{corr sales}_{j,b,t}}{\sum_{j' \in \text{NACE 4d}(j)} \text{corr sales}_{j',b,t}}.
\]
This is the buyer's share of cell-level spending going to a supplier whose costs are exposed to ETS at intensity $\text{firm cost share}_{j,t}$. It answers: ``within the NACE 4-digit category, how much of the buyer's spending in that category is exposed to the ETS shock through this specific supplier?''

\paragraph{Pair-shock at the total-cost denominator.} The same numerator denominated at the buyer's total-input level:
\[
    \text{pair-shock-total}_{j,b,t} \;=\; \text{firm cost share}_{j,t} \;\times\; \frac{\text{corr sales}_{j,b,t}}{\text{total inputs}_{b,t}}.
\]
This answers: ``what fraction of the buyer's total input bill is at risk under full pass-through if this specific supplier passes its carbon cost forward?''

\paragraph{Buyer-total-shock.} The sum of pair-shock-total across all the buyer's ETS suppliers in the year:
\[
    \text{buyer-total-shock}_{b,t} \;=\; \sum_{j \in \text{ETS}} \text{pair-shock-total}_{j,b,t}.
\]
This is the buyer's aggregate exposure --- the total fraction of input cost that would absorb the EUA shock at full pass-through if the buyer reallocated nothing.

Phase IV distributions are reported in Table~\ref{tab:buyer_shock_phase4}.

\begin{table}[H]
    \centering
    \caption{Buyer-side shock magnitudes in Phase~IV (2021--22), at three denominator scales.}
    \label{tab:buyer_shock_phase4}
    \footnotesize
    \begin{tabular}{lrrrrrr}
        \toprule
        Quantity & p50 & p90 & p95 & p99 & mean & sales-wmean \\
        \midrule
        Pair-shock (NACE-4d denom., $n_{\text{pairs}} = 3{,}597$)   & $0.084\%$ & $3.30\%$ & $8.29\%$  & $12.0\%$ & $1.22\%$ & $0.94\%$ \\
        Pair-shock-total (total-input denom., $n_{\text{pairs}} = 3{,}541$)   & $0.0016\%$ & $1.16\%$ & --- & $\mathbf{39.1\%}$  & $1.56\%$ & $6.74\%$ \\
        Buyer-total-shock ($n_{\text{buyers}} = 15{,}442$)         & $0$       & $0$      & $0.031\%$ & $6.90\%$ & $0.36\%$ & --- \\
        \bottomrule
    \end{tabular}
\end{table}

Two readings of Table~\ref{tab:buyer_shock_phase4} are informative.

First, \emph{the typical buyer faces no carbon-cost shock at all.} Buyer-total-shock is exactly zero at p50 and p90: most buyers in the regulated-intensive sample do not have any ETS suppliers in their portfolio. The right tail is non-trivial (p99 = 6.9\%) and concentrated in a small number of buyer sectors.

Second, \emph{at the cell where the substitution test has economic content, the right tail is large but the median is small.} Pair-shock at the NACE-4d denominator has p50 of $0.08\%$, p90 of $3.30\%$, and p99 of $12.0\%$. The p99 cell --- where a 99th-percentile-exposed buyer faces a $\sim$12\% pair-shock at the cell level --- is at the same scale as the per-input tariff cuts (~30 percentage points) that \citet{peter2023} use to identify $\sigma \approx 4$ on Indian intermediates. At the buyer's total-cost level the right tail is even more concentrated, with pair-shock-total p99 reaching $39.1\%$ across the population.

\subsection{Comparison to the buyer-side noise floor}\label{sec:shock_signal_to_noise}

The relevant question for the leakage tests is not whether the shock is large in absolute terms, but whether it is large relative to the natural noise in the buyer's input-cost structure --- the cost-share variation that buyers absorb without observable reallocation under non-ETS conditions. We measure this noise floor as
\[
    \sigma_{\text{share}} \;=\; \text{within-firm sd of } \Delta \log(\text{inputs}_{b,t}/\text{revenue}_{b,t})
\]
on the 2005--2019 (pre-Phase-IV) window, scale-stripped to isolate the price-and-mix component. Population p50 of $\sigma_{\text{share}}$ is $15.0\%$.

The signal-to-noise comparison at the matched scale (pair-shock-total vs $\sigma_{\text{share}}$) is reported in Table~\ref{tab:signal_to_noise}.

\begin{table}[H]
    \centering
    \caption{Signal-to-noise: Phase~IV pair-shock-total p90 relative to pre-shock $\sigma_{\text{share}}$ p50.}
    \label{tab:signal_to_noise}
    \footnotesize
    \begin{tabular}{lrrr}
        \toprule
        Subset & Pair-shock-total p90 & $\sigma_{\text{share}}$ p50 & Signal-to-noise \\
        \midrule
        Population (all treated active pairs) & $1.16\%$ & $15.0\%$ & $0.08\sigma$ \\
        \midrule
        \emph{Comparison: \citet{heise2024} 1-$\sigma$ FX shock (annualized, buyer-side)} & $\sim 1.5\%$ & --- & $\sim 0.10\sigma$ \\
        \bottomrule
    \end{tabular}
\end{table}

\textbf{Across regulated sectors, the typical pair-shock is well below the buyer-side input-cost noise floor.} Population pair-shock-total at $0.08\sigma$ of $\sigma_{\text{share}}$ is below the noise floor by a factor of 12. ``No substitution'' at the typical buyer is what we would expect even in a model with substantial $\sigma$, because the signal is too small for buyers to act on specifically.

